\documentclass[11pt]{article}

\usepackage[T1]{fontenc}
\usepackage[utf8]{inputenc}
\usepackage{amsmath,amssymb,amsthm,mathtools}
\usepackage{microtype}
\usepackage{hyperref}
\usepackage{geometry}
\geometry{margin=1in}

\newtheorem{theorem}{Theorem}[section]
\newtheorem{proposition}[theorem]{Proposition}
\newtheorem{lemma}[theorem]{Lemma}
\newtheorem{corollary}[theorem]{Corollary}
\newtheorem{remark}[theorem]{Remark}
\newtheorem{definition}[theorem]{Definition}

\newcommand{\C}{\mathbb{C}}
\newcommand{\op}{\mathrm{op}}
\newcommand{\cmax}{C_{\max}}
\newcommand{\dd}{\delta}
\newcommand{\etaq}{\eta}
\newcommand{\tauq}{\tau}
\newcommand{\PhiRho}{\Phi_{\rho}}

\title{Sharp Quantitative Stability of Finite Projective Records\\
under Approximately Commuting Unitary Transformations}

\author{Bertrand Dalimier\\
Independent Researcher}

\date{}

\begin{document}

\maketitle

\begin{abstract}
Let $(P_i)_{i\in I}$ be a finite orthogonal projective decomposition of a
finite-dimensional complex Hilbert space and let $U$ be unitary.  Write
$a_i=\|[P_i,U]\|_{\op}$ for the cellwise commutator defects,
$\dd^2=\sum_i a_i^2$ for their global quadratic budget, and
\[
  \cmax=\max_{S\subseteq I}\|[P_S,U]\|_{\op},
  \qquad P_S=\sum_{i\in S}P_i ,
\]
for the finite cut envelope.  We prove the sharp universal estimate
$\cmax^2\le \dd^2/2$.  When $\dd>0$, equality holds if and only if exactly two
cells have nonzero commutator defect.  We then establish an inverse stability
theorem: if
\[
  \cmax^2\ge (1-\etaq)\frac{\dd^2}{2},\qquad
  0\le\etaq\le\frac15,\qquad \dd\ge\rho>0,
\]
then two cells capture all but a relative fraction
\[
  \PhiRho(\etaq)
  =2\etaq+\frac{36\etaq}{(1-3\etaq)\rho^2}
\]
of the global quadratic budget.  For $\etaq\le1/6$ this is bounded by
$\etaq(2+72/\rho^2)$.  Finally, an explicit three-cell family shows that the
inverse-square scale dependence is necessary up to constants: in an iterated
near-saturation regime, $(\tauq/\etaq)\dd^2\to4$.  All principal statements
and the sharpness construction have been formally verified in Lean 4.
\end{abstract}

\noindent\textbf{Keywords.}
projector commutators; quantitative stability; rigidity; near equality;
unitary operators; formal verification

\noindent\textbf{2020 Mathematics Subject Classification.}
15A60 (primary); 15A45, 47A63 (secondary).

\section{Introduction}

Finite projective decompositions provide a simple setting in which a unitary
transformation can be compared with a prescribed collection of mutually
orthogonal sectors.  If $P_i$ is the projector associated with sector $i$,
then the commutator $[P_i,U]$ measures the failure of $U$ to preserve that
sector.  The individual operator norms
\[
  a_i:=\|[P_i,U]\|_{\op}
\]
therefore define a nonnegative finite defect profile.  The quadratic aggregate
\[
  \dd^2:=\sum_{i\in I} a_i^2
\]
is a natural dimension-free global measure of failure of simultaneous
commutation.

A second object is obtained by grouping sectors into a cut.  For
$S\subseteq I$, let
\[
  P_S:=\sum_{i\in S}P_i,\qquad
  C_S:=\|[P_S,U]\|_{\op}.
\]
Since $I$ is finite, the cut envelope
\[
  \cmax:=\max_{S\subseteq I}C_S
\]
is attained.  The basic inequality studied here is
\begin{equation}\label{eq:envelope}
  \cmax^2\le \frac{\dd^2}{2}.
\end{equation}
It is a finite operator-theoretic inequality: no probability assignment,
dynamical assumption, open-system structure, or asymptotic Hilbert-space
limit is involved.

The purpose of this paper is to analyze the near-equality regime of
\eqref{eq:envelope}.  There are three levels.

First, we identify the exact equality cases.  For positive global defect,
saturation of \eqref{eq:envelope} is equivalent to the statement that exactly
two cells are active: all other cellwise commutator defects vanish.

Second, we prove a quantitative inverse theorem.  If the cut envelope nearly
saturates \eqref{eq:envelope}, then almost all of the quadratic defect budget
must be carried by two cells.  The estimate is independent of the number of
cells.  At a fixed positive defect scale $\rho$, the relative two-cell tail is
$O(\etaq/\rho^2)$ as the saturation deficit $\etaq$ tends to zero.

Third, we construct an explicit three-cell family showing that the
$1/\rho^2$ scale is not a proof artifact.  The family has closed formulas for
the three cell budgets, the global defect, the maximal cut, the
near-saturation parameter, and the two-cell tail.  Its asymptotics give the
constant
\[
  \lim_{m\to\infty}\left[
    \lim_{n\to\infty}\frac{\tauq_{n,m}}{\etaq_{n,m}}\dd_{n,m}^2
  \right]=4.
\]
Thus a stability estimate uniform down to zero defect cannot in general have
a tail bound of order $O(\etaq)$ with a scale-independent constant.

The arguments separate into a finite real concentration layer and an
operator-theoretic layer.  This separation is also reflected in a companion
Lean 4 formalization.  The principal theorem statements, equality
classification, rank-one formula, explicit sharpness family, and asymptotic
limit are machine checked.  The formal development is used here as a
verification artifact rather than as the mathematical subject of the paper.

% TODO before submission:
% 1. Complete systematic priority/related-work audit.
% 2. Replace repository commit by archival release/DOI.
% 3. Adapt notation/template to target journal.

\section{Finite projective setup}

Let $H$ be a finite-dimensional complex Hilbert space.  Let
$I$ be a finite index set and let $(P_i)_{i\in I}$ be orthogonal projections
satisfying
\begin{equation}\label{eq:decomp}
  P_iP_j=0\quad(i\neq j),
  \qquad
  \sum_{i\in I}P_i=I_H.
\end{equation}
We call the family a finite projective record or finite projective
decomposition.  Let $U:H\to H$ be unitary.

For each cell $i$, define
\begin{equation}\label{eq:ai}
  a_i:=\|[P_i,U]\|_{\op},
  \qquad
  p_i:=a_i^2,
\end{equation}
and define the global quadratic commutator defect
\begin{equation}\label{eq:delta}
  \dd:=\left(\sum_{i\in I}p_i\right)^{1/2}.
\end{equation}
For a cut $S\subseteq I$, set
\begin{equation}\label{eq:cut}
  P_S:=\sum_{i\in S}P_i,
  \qquad
  C_S:=\|[P_S,U]\|_{\op},
  \qquad
  \cmax:=\max_{S\subseteq I}C_S.
\end{equation}

The quantity $\dd$ is the $\ell^2$ norm of the cellwise operator-norm profile.
The quantity $\cmax$ instead probes the largest coherent defect visible after
aggregating an arbitrary subset of the cells.

For later use, when $\dd>0$ define the optimal two-cell tail
\begin{equation}\label{eq:tau}
  \tauq:=
  \min_{\substack{i,j\in I\\i\neq j}}
  \frac{\dd^2-p_i-p_j}{\dd^2}.
\end{equation}
Equivalently, $\tauq$ is the fraction of the total quadratic budget left after
retaining the best pair of distinct cells.

\begin{remark}
The formal theorem is stated without division by $\dd^2$: it asserts the
existence of distinct $i,j$ such that
\[
  \dd^2-p_i-p_j\le \varepsilon\,\dd^2.
\]
This formulation remains meaningful at zero defect.  In the positive-defect
regime it is equivalent to $\tauq\le\varepsilon$.
\end{remark}

\section{The finite cut envelope}

The key observation is that a cut commutator is controlled simultaneously by
the defect budget on the cut and by the defect budget on its complement.

Relative to the orthogonal splitting
\[
  H=P_SH\oplus (I-P_S)H,
\]
the commutator $[P_S,U]$ is off diagonal.  Its two cross blocks are
\[
  P_SU(I-P_S),
  \qquad
  (I-P_S)UP_S.
\]
For a unitary $U$ the operator norm of the cut commutator is the maximum of
the norms of these two cross blocks.  The incoming block is controlled by the
cellwise budgets on $S$, while the outgoing block is controlled by the
budgets on $S^c$.  Consequently,
\begin{equation}\label{eq:cut-min}
  C_S^2
  \le
  \min\left\{
    \sum_{i\in S}p_i,\,
    \sum_{i\notin S}p_i
  \right\}.
\end{equation}

Since the two quantities inside the minimum sum to $\dd^2$, we obtain the
universal envelope.

\begin{theorem}[Cut-envelope inequality]\label{thm:envelope}
For every finite projective decomposition and every unitary $U$,
\[
  \boxed{\cmax^2\le\frac{\dd^2}{2}}.
\]
\end{theorem}

\begin{proof}
For every $S\subseteq I$, \eqref{eq:cut-min} gives
\[
  C_S^2
  \le
  \min\{A_S,\dd^2-A_S\},
  \qquad
  A_S:=\sum_{i\in S}p_i.
\]
The elementary inequality
$\min\{x,T-x\}\le T/2$ for $0\le x\le T$ yields
$C_S^2\le\dd^2/2$.  Taking the maximum over the finite set of cuts proves the
claim.
\end{proof}

The coefficient $1/2$ is optimal.  More importantly for the present paper,
its equality cases are rigid.

\section{Quantitative stability of near-saturators}

We now prove the inverse theorem behind Theorem~\ref{thm:rigidity}.  For a
cut $S\subseteq I$, write
\[
  A_S:=\sum_{i\in S}p_i,
  \qquad
  C_S:=\|[P_S,U]\|_{\op}.
\]

\subsection{A one-sided concentration lemma}

The numerical constant in the final stability theorem enters through the
following dimension-free lemma.

\begin{lemma}[One-sided concentration]\label{lem:onesided}
Let $S$ be a cut with $C_S>0$, and put
\[
  e_S:=A_S-C_S^2.
\]
If $e_S\le C_S^2/2$, then there exists $i\in S$ such that
\begin{equation}\label{eq:onesided}
  A_S-p_i
  \le
  2e_S+\frac{18e_S}{C_S^2-e_S}.
\end{equation}
\end{lemma}

\begin{proof}
Relative to $P_SH\oplus P_{S^c}H$, the cut commutator has incoming and
outgoing cross blocks.  We give the argument when the incoming block realizes
$C_S$; the outgoing case follows by the same argument after exchanging the
two blocks.

Because the space is finite dimensional, choose a unit vector
$v\in P_{S^c}H$ on which the incoming block attains its norm.  For $c\in S$
set
\[
  z_c:=U^*P_cUv,\qquad
  r_c:=\|z_c\|^2,\qquad
  \alpha_c:=p_c-r_c.
\]
The vectors $z_c$ are mutually orthogonal and the witness identities give
\begin{equation}\label{eq:witness-defect}
  \alpha_c\ge0,\qquad
  \sum_{c\in S}\alpha_c=e_S.
\end{equation}
Call a cell good when $\alpha_c\le p_c/2$.  The total $p$-mass of the bad
cells is at most
\begin{equation}\label{eq:badmass}
  \sum_{\text{bad }c}p_c\le2e_S.
\end{equation}
Discarding good cells with $p_c=0$ does not change the good-cell mass.  On
the remaining active good cells one has $r_c>0$.  Define
\[
  \xi_c:=\frac{z_c}{\sqrt{r_c}}
\]
and a normalized comparison family $\zeta_c$ obtained by replacing the
residual component of $\xi_c$ by $\sqrt{p_c}\,v$.  The witness geometry gives
three facts:
\begin{align}
  &(\xi_c)_c\ \text{is orthonormal}, \label{eq:xi-orth}\\
  &p_cp_d=|\langle\zeta_c,\zeta_d\rangle|^2
      \quad(c\ne d), \label{eq:zeta-gram}\\
  &\sum_c\|\zeta_c-\xi_c\|^2\le2e_S. \label{eq:gram-close}
\end{align}

For completeness, if $\xi$ is any orthonormal family and $\zeta$ satisfies
\eqref{eq:zeta-gram}, with
$q:=\sum_c\|\zeta_c-\xi_c\|^2$, expansion of each off-diagonal inner product
around the orthonormal reference family and Bessel's inequality give
\begin{equation}\label{eq:gram-spread-general}
  \sum_c\sum_{d\ne c}p_cp_d
  \le 6q+3q^2.
\end{equation}
Here $C_S\le1$, hence $e_S\le1/2$.  By \eqref{eq:gram-close},
$q\le2e_S$, so
\begin{equation}\label{eq:gram18}
  \sum_c\sum_{d\ne c}p_cp_d\le18e_S.
\end{equation}
For a finite nonnegative profile,
\[
  \sum_c\sum_{d\ne c}p_cp_d
  =
  \left(\sum_cp_c\right)^2-\sum_cp_c^2.
\]
Choosing a maximal entry therefore shows that if the total active-good mass
is at least $L>0$, one active good cell leaves active-good tail at most
$18e_S/L$.

The witness identities and \eqref{eq:witness-defect} give the lower bound
\[
  L=C_S^2-e_S
\]
for that active-good mass.  Adding the bad-cell contribution
\eqref{eq:badmass} yields \eqref{eq:onesided}.
\end{proof}

\subsection{Bilateral concentration across a cut}

Let
\[
  A:=\dd^2,\qquad C:=C_S,\qquad
  A_T:=A_{S^c},
\]
and define
\[
  e_S:=A_S-C^2,\qquad
  e_T:=A_T-C^2,\qquad
  E:=A-2C^2.
\]
The complementary budget identity $A_S+A_T=A$ gives
\begin{equation}\label{eq:e-split}
  e_S+e_T=E.
\end{equation}

\begin{proposition}[Bilateral cut concentration]\label{prop:bilateral}
Assume $C>0$ and $E\le C^2/2$.  Then there exist
$i\in S$ and $j\in S^c$ such that
\begin{equation}\label{eq:bilateral}
  A-p_i-p_j
  \le
  2E+\frac{18E}{C^2-E}.
\end{equation}
\end{proposition}

\begin{proof}
The cut estimate gives $C^2\le A_S$ and $C^2\le A_T$, hence
$e_S,e_T\ge0$.  By \eqref{eq:e-split}, each is at most $E$ and therefore at
most $C^2/2$.  Lemma~\ref{lem:onesided}, applied to $S$ and $S^c$, gives
$i\in S$ and $j\in S^c$ such that
\[
  A_S-p_i
  \le
  2e_S+\frac{18e_S}{C^2-e_S},
\]
\[
  A_T-p_j
  \le
  2e_T+\frac{18e_T}{C^2-e_T}.
\]
Set $d:=C^2-E>0$.  Since
$d\le C^2-e_S$ and $d\le C^2-e_T$,
\[
  \frac{e_S}{C^2-e_S}
  +
  \frac{e_T}{C^2-e_T}
  \le
  \frac{e_S+e_T}{d}
  =
  \frac Ed.
\]
Adding the two one-sided inequalities proves
\eqref{eq:bilateral}.  Since $i\in S$ and $j\in S^c$, the two cells are
distinct.
\end{proof}

\subsection{The scale-sensitive modulus}

Let $0\le\etaq\le1/5$ and assume
\begin{equation}\label{eq:near}
  \cmax^2
  \ge
  (1-\etaq)\frac{\dd^2}{2}.
\end{equation}
We also impose the positive scale condition
\begin{equation}\label{eq:scale}
  \dd\ge\rho>0.
\end{equation}
Define
\begin{equation}\label{eq:modulus}
  \PhiRho(\etaq):=
  2\etaq+
  \frac{36\etaq}{(1-3\etaq)\rho^2}.
\end{equation}

\begin{theorem}[Quantitative two-cell stability]\label{thm:stability}
Suppose \eqref{eq:near} and \eqref{eq:scale} hold with
$0\le\etaq\le1/5$.  Then there exist distinct cells $i,j$ such that
\begin{equation}\label{eq:stability-conclusion}
  \boxed{
  \dd^2-p_i-p_j
  \le
  \PhiRho(\etaq)\,\dd^2.
  }
\end{equation}
Equivalently,
\[
  \tauq\le
  2\etaq+
  \frac{36\etaq}{(1-3\etaq)\rho^2}.
\]
\end{theorem}

\begin{proof}
Choose a cut $S$ attaining $\cmax$ and put
\[
  A:=\dd^2,\qquad C:=\cmax,\qquad E:=A-2C^2.
\]
Theorem~\ref{thm:envelope} gives $E\ge0$, while
\eqref{eq:near} gives
\begin{equation}\label{eq:Eeta}
  E\le\etaq A.
\end{equation}
Since $\etaq\le1/5$ and $A=2C^2+E$,
\[
  E\le\frac A5
  \quad\Longrightarrow\quad
  E\le\frac{C^2}{2}.
\]
Moreover $A>0$, and the same inequalities imply $C>0$.  Proposition
\ref{prop:bilateral} therefore gives distinct $i,j$ with
\[
  A-p_i-p_j
  \le
  2E+\frac{18E}{d},
  \qquad
  d:=C^2-E.
\]
Because $C^2=(A-E)/2$,
\begin{equation}\label{eq:didentity}
  d=\frac{A-3E}{2}.
\end{equation}
Using \eqref{eq:Eeta},
\[
  d\ge\frac{(1-3\etaq)A}{2}>0
\]
and hence
\[
  \frac Ed
  \le
  \frac{2\etaq}{1-3\etaq}.
\]
Thus
\begin{equation}\label{eq:absolute-tail}
  A-p_i-p_j
  \le
  2\etaq A+
  \frac{36\etaq}{1-3\etaq}.
\end{equation}
Finally $A=\dd^2\ge\rho^2$, so
\[
  \frac{36\etaq}{1-3\etaq}
  \le
  A\,
  \frac{36\etaq}{(1-3\etaq)\rho^2}.
\]
Substitution into \eqref{eq:absolute-tail} gives
\eqref{eq:stability-conclusion}.
\end{proof}

At every fixed positive scale,
\[
  \PhiRho(\etaq)\longrightarrow0
  \qquad(\etaq\downarrow0).
\]

\begin{corollary}[Transparent linear form]\label{cor:linear}
If in addition $\etaq\le1/6$, then
\begin{equation}\label{eq:linear}
  \boxed{
    \tauq\le
    \etaq\left(2+\frac{72}{\rho^2}\right).
  }
\end{equation}
\end{corollary}

\begin{proof}
For $\etaq\le1/6$ one has
$1-3\etaq\ge1/2$, so
\[
  \frac{36\etaq}{(1-3\etaq)\rho^2}
  \le
  \frac{72\etaq}{\rho^2}.
\]
Apply Theorem~\ref{thm:stability}.
\end{proof}


\section{Exact rigidity at saturation}

Call a cell \emph{active} when $p_i>0$.

\begin{theorem}[Exact equality classification]\label{thm:rigidity}
Assume $\dd>0$.  Then
\[
  \boxed{
    \cmax^2=\frac{\dd^2}{2}
    \quad\Longleftrightarrow\quad
    \text{exactly two cells are active}.
  }
\]
More explicitly, equality holds if and only if there exist distinct
$i,j\in I$ such that
\[
  p_i>0,\qquad p_j>0,\qquad
  p_k=0\quad\text{for every }k\notin\{i,j\}.
\]
\end{theorem}

\begin{proof}
Set $A=\dd^2$.  First assume saturation,
$\cmax^2=A/2$.  Apply Theorem~\ref{thm:stability} with
$\etaq=0$ and $\rho=\dd>0$.  Since $\Phi_{\dd}(0)=0$, there are distinct
$i,j$ such that
\[
  A-p_i-p_j\le0.
\]
On the other hand,
\[
  A-p_i-p_j=\sum_{k\notin\{i,j\}}p_k\ge0,
\]
so every $p_k$ with $k\notin\{i,j\}$ vanishes.

It remains to show that the two selected cells are genuinely active.  A
singleton is an admissible cut, and therefore
\[
  p_i=\|[P_i,U]\|_{\op}^2\le\cmax^2=\frac A2,
  \qquad
  p_j\le\frac A2.
\]
Since the remaining cells carry zero budget, $p_i+p_j=A$.  Hence
$p_i=p_j=A/2>0$.

Conversely, suppose exactly two cells $i\ne j$ are active.  Then
$A=p_i+p_j$, so
\[
  \max\{p_i,p_j\}\ge\frac A2.
\]
Again using singleton cuts,
\[
  \cmax^2\ge\max\{p_i,p_j\}\ge\frac A2.
\]
The cut-envelope inequality, Theorem~\ref{thm:envelope}, gives the reverse
bound $\cmax^2\le A/2$, and equality follows.
\end{proof}

The hypothesis $\dd>0$ only removes the degenerate configuration in which
every commutator vanishes.


\section{An explicit three-cell sharpness family}\label{sec:sharp}

We construct a family in dimension three whose near-saturation deficit tends
to zero while the ratio $\tauq/\etaq$ diverges at the inverse-square defect
scale.

Let $(p_n,q_n)$ be the positive integer recurrence
\[
  (p_0,q_0)=(1,1),
\]
\[
  p_{n+1}=3p_n+4q_n,
  \qquad
  q_{n+1}=2p_n+3q_n.
\]
It preserves
\begin{equation}\label{eq:pell}
  p_n^2+1=2q_n^2.
\end{equation}
Define
\[
  a_n=\frac{p_n+1}{2q_n},
  \qquad
  b_n=\frac{p_n-1}{2q_n},
\]
so that $a_n^2+b_n^2=1$.  Rotate these amplitudes by $45^\circ$:
\begin{equation}\label{eq:AB}
  A_n=\frac{a_n+b_n}{\sqrt2},
  \qquad
  B_n=\frac{a_n-b_n}{\sqrt2}.
\end{equation}
Then
\[
  A_n^2+B_n^2=1,\qquad
  B_n=\frac{1}{\sqrt2\,q_n},
\]
hence
\begin{equation}\label{eq:AB-limits}
  B_n^2\to0,\qquad A_n^2\to1.
\end{equation}

For a second index $m$, let $c_m,s_m$ be the explicit rational unit-circle
parameters
\[
  c_m^2+s_m^2=1,\qquad s_m>0,
\]
coming from the same Pell family, and set
\begin{equation}\label{eq:d}
  d_m:=1-c_m.
\end{equation}
Thus
\begin{equation}\label{eq:s-d}
  s_m^2=2d_m-d_m^2,
  \qquad
  d_m\to0.
\end{equation}

The three-dimensional unitary is obtained by conjugating the planar rotation
with a basis change whose first direction is
$(A_n,B_n,0)$.  For the three rank-one cells, the exact squared commutator
budgets are
\begin{align}
  p_0 &=
  2d_mA_n^2-d_m^2A_n^4, \label{eq:p0}\\
  p_1 &=
  2d_mB_n^2-d_m^2B_n^4, \label{eq:p1}\\
  p_2 &= s_m^2. \label{eq:p2}
\end{align}
These identities follow from the rank-one formula
\begin{equation}\label{eq:rankone}
  \boxed{
  \|[P_e,U]\|_{\op}^2
  =
  1-\bigl|\langle e,Ue\rangle\bigr|^2
  }
\end{equation}
for a unit vector $e$ and the rank-one projector $P_e$.

Summing \eqref{eq:p0}--\eqref{eq:p2} gives
\begin{equation}\label{eq:global-sharp}
  \boxed{
  \dd_{n,m}^2
  =
  2s_m^2+
  2d_m^2A_n^2B_n^2.
  }
\end{equation}
The maximal cut is exact:
\begin{equation}\label{eq:cmax-sharp}
  \boxed{
  \cmax(n,m)=s_m.
  }
\end{equation}

Define the relative saturation deficit
\[
  \etaq_{n,m}
  :=
  \frac{\dd_{n,m}^2-2\cmax(n,m)^2}{\dd_{n,m}^2}.
\]
Then
\begin{equation}\label{eq:eta-sharp}
  \boxed{
  \etaq_{n,m}
  =
  \frac{
    2d_m^2A_n^2B_n^2
  }{
    2s_m^2+2d_m^2A_n^2B_n^2
  }.
  }
\end{equation}
Retaining cells $0$ and $2$ leaves the cell-$1$ tail
\begin{equation}\label{eq:tau-sharp}
  \tauq_{n,m}
  =
  \frac{
    2d_mB_n^2-d_m^2B_n^4
  }{
    2s_m^2+2d_m^2A_n^2B_n^2
  }.
\end{equation}
The exact quotient simplifies to
\begin{equation}\label{eq:ratio-sharp}
  \boxed{
  \frac{\tauq_{n,m}}{\etaq_{n,m}}
  =
  \frac{2-d_mB_n^2}{2d_mA_n^2}.
  }
\end{equation}

\begin{theorem}[Scale sharpness]\label{thm:scale-sharp}
For fixed $m$,
\begin{equation}\label{eq:fixed-m-ratio}
  \lim_{n\to\infty}
  \frac{\tauq_{n,m}}{\etaq_{n,m}}
  =
  \frac1{d_m},
\end{equation}
and
\begin{equation}\label{eq:fixed-m-delta}
  \lim_{n\to\infty}\dd_{n,m}^2=2s_m^2.
\end{equation}
Consequently,
\begin{equation}\label{eq:iterated4}
  \boxed{
  \lim_{m\to\infty}
  \left[
    \lim_{n\to\infty}
    \frac{\tauq_{n,m}}{\etaq_{n,m}}
    \dd_{n,m}^2
  \right]
  =4.
  }
\end{equation}
\end{theorem}

\begin{proof}
Equations \eqref{eq:fixed-m-ratio} and \eqref{eq:fixed-m-delta} follow from
\eqref{eq:AB-limits}, \eqref{eq:global-sharp}, and
\eqref{eq:ratio-sharp}.  Their product tends, for fixed $m$, to
\[
  \frac{2s_m^2}{d_m}.
\]
Using \eqref{eq:s-d},
\[
  \frac{2s_m^2}{d_m}
  =
  4-2d_m.
\]
Finally $d_m\to0$, giving \eqref{eq:iterated4}.
\end{proof}

\begin{corollary}[Necessity of the inverse-square scale]
Along the family above,
\[
  \frac{\tauq_{n,m}}{\etaq_{n,m}}
  \asymp \frac{1}{\dd_{n,m}^2}
\]
in the iterated small-defect regime.  Hence a dimension-free near-equality
theorem with tail $O(\etaq)$ and a constant independent of the defect scale
cannot hold in general.  The power $1/\rho^2$ in
Theorem~\ref{thm:stability} is therefore necessary up to multiplicative
constants.
\end{corollary}

\section{Relation with the sharp record-transfer bound}

The cut envelope also enters the sharp uniform transfer inequality for finite
projective record profiles.  For a normalized state $x$, the companion
formal development proves
\begin{equation}\label{eq:transfer}
  \|r_D(Ux)-r_D(x)\|_1
  \le
  \sqrt2\,\dd,
\end{equation}
with the universal coefficient $\sqrt2$ optimal.

The present paper addresses the inverse geometry behind near saturation of
the cut stage that feeds \eqref{eq:transfer}.  Theorems
\ref{thm:rigidity} and \ref{thm:stability} show that extremal and
near-extremal cut behavior is effectively two-cell.  This part of the
argument is independent of the Born-rule interpretation of the record
profile: it depends only on finite projectors, unitary maps, operator norms,
and finite real concentration.

A possible extension is to combine the present cut-envelope stability with a
separate stability analysis of every remaining inequality in the derivation
of \eqref{eq:transfer}.  Such a result would classify near-extremizers of the
full $\sqrt2$ transfer inequality, rather than only the operator-theoretic cut
stage.  We do not claim such a classification here.

\section{Formal verification}

The mathematical core has a companion Lean 4 formalization in the repository
\[
  \texttt{Bobart0/everettian-decoherence-lean}.
\]
The development is organized so that the publication-facing surface imports
the exact rigidity theorem, the quantitative stability theorem, and the
sharpness asymptotics.  A dedicated audit module checks the principal
declarations and prints their axiomatic dependencies.

At manuscript version 0.1, the reference development is the CI-verified
repository state containing the following principal declarations:
\begin{center}
\begin{tabular}{p{0.39\textwidth}p{0.52\textwidth}}
\hline
Mathematical role & Lean declaration\\
\hline
Cut envelope &
\texttt{maxSubsetCommutatorOpNorm\_sq\_le\_half\_globalSq}\\
Exact rigidity &
\texttt{maxCut\_saturation\_iff\_exactlyTwoActiveCellCommutators}\\
Quantitative modulus &
\texttt{twoCellCommutatorConcentratedWithin\_of\_maxCut\_near\_saturated}\\
Linear corollary &
\texttt{twoCellCommutatorConcentratedWithin\_linear\_of\_maxCut\_near\_saturated}\\
Rank-one identity &
\texttt{perspectiveProjectorCommutatorOpNormProfile\_sq\_rankOne}\\
Sharpness ratio &
\texttt{stabilitySharpnessTau\_div\_Eta\_exact}\\
Fixed-scale asymptotics &
\texttt{stabilitySharpnessTau\_div\_Eta\_tendsto\_inv\_D}\\
Defect-scale asymptotics &
\texttt{stabilitySharpness\_delta\_sq\_tendsto\_two\_s\_sq}\\
Constant-$4$ obstruction &
\texttt{stabilitySharpnessScaleLimit\_tendsto\_four}\\
\hline
\end{tabular}
\end{center}

The publication-facing facade is
\texttt{EverettianDecoherence.Approximation.StabilitySharpness}; the dedicated
audit module is
\texttt{EverettianDecoherence.Audit.StabilityRigiditySharpness}.
The final formalization head preceding this manuscript branch is
\texttt{81eb3ff53f6219749b84491513ff437fe7c545c1}, whose complete CI passed,
including the library build, audit aggregate, guards, and diff-hygiene checks.

For archival submission, this commit reference should be replaced or
supplemented by a tagged release and a persistent archive identifier.

\paragraph{AI-assisted formalization workflow.}
OpenAI ChatGPT was used as an interactive assistant during parts of the Lean
development, including proof-structure exploration, code drafting, and
diagnosis of compiler errors. Every publication-facing declaration reported
here is checked by Lean itself, and the repository's continuous-integration
workflow rebuilds the library and the dedicated audit surface. The author
reviewed the mathematical statements and proof dependencies and remains
responsible for the correctness and interpretation of the results.

\section{Discussion and limitations}

The stability theorem is dimension-free in the number of projective cells:
its constants do not grow with $|I|$.  The price is the explicit positive
scale $\rho$.  The sharpness family shows that this dependence is structural
rather than merely technical.

The theorem is also deliberately algebraic.  It does not assert that the
projective decomposition is dynamically selected, that the sectors are
environmentally redundant, that off-diagonal density-matrix elements decay
in time, or that a preferred basis emerges.  Nor does the operator-theoretic
core require a probability interpretation of the cells.  Such physical
questions require additional hypotheses that are outside the scope of the
present result.

The constants $36$ and $72$ arise from the quantitative proof architecture
and are not claimed to be optimal.  The sharpness construction identifies the
correct inverse-square scale and an asymptotic constant $4$ for the explicit
family, but it does not by itself determine the best universal coefficient in
the full finite-dimensional stability theorem.

\section{Relation to existing operator and matrix theory}
\label{sec:related}

The closest classical background is the geometry of projections and invariant
subspaces.  Halmos' two-subspaces theorem gives a canonical block
representation for a pair of closed subspaces, equivalently for a pair of
orthogonal projections \cite{Halmos1969TwoSubspaces}.  Davis--Kahan
perturbation theory quantifies rotation of invariant subspaces using canonical
angles \cite{DavisKahan1970}.  These results concern the relative geometry or
perturbation of one or two subspaces.  Here the datum is instead a complete
finite orthogonal decomposition $(P_i)_{i\in I}$ together with a unitary $U$,
and the object studied is the entire cellwise profile
$\{\|[P_i,U]\|_{\op}\}_{i\in I}$ together with the maximum over all aggregate
cuts.

Norm inequalities for commutators form a second neighboring literature.
Bhatia, Davis, and Kittaneh proved general inequalities for commutators in
unitarily invariant norms with applications to spectral variation
\cite{BhatiaDavisKittaneh1991}; further operator-norm inequalities were
developed by Kittaneh \cite{Kittaneh2007}.  Structural descriptions are also
known for commutators of pairs of orthogonal projections
\cite{Li2004}.  Those works provide context for individual commutators, while
our central finite-family functional is
\[
  \max_{S\subseteq I}
  \left\|
    \left[\sum_{i\in S}P_i,U\right]
  \right\|_{\op},
\]
compared to the quadratic budget
$\sum_i\|[P_i,U]\|_{\op}^2$.

A third nearby question asks whether almost commuting operators can be
perturbed to exactly commuting ones.  The answer depends strongly on the
operator class.  Voiculescu exhibited asymptotically commuting finite-rank
unitaries without commuting unitary approximants
\cite{Voiculescu1983}; Exel and Loring developed the corresponding
topological obstruction theory \cite{ExelLoring1989,ExelLoring1991}.  For
self-adjoint matrices, Lin proved dimension-independent approximation by
commuting self-adjoints \cite{Lin1997AlmostCommuting}; see also
\cite{FriisRordam1996,Hastings2009}.  Effective bounds for almost commuting
unitaries remain an active topic \cite{DorOnHallKachkovskiy2025}, and recent
work formulates projection characterizations for operators admitting nearby
commuting approximants \cite{Herrera2024,Herrera2026}.  Almost-invariant
subspaces and approximate commutation provide another related perspective
\cite{MarcouxPopovRadjavi2013}.

The problem considered here is different from the commuting-approximant
problem.  The projectors $(P_i)$ are fixed and mutually orthogonal, and the
unitary $U$ is not replaced by another operator.  Instead, we ask what the
\emph{existing commutator profile} must look like when one aggregate cut
nearly attains its universal envelope.  The resulting rigidity is therefore
a near-equality statement internal to a fixed finite decomposition:
saturation forces exactly two active cells, near saturation forces
quantitative two-cell concentration, and the three-cell construction shows
that the inverse-square defect scale is unavoidable up to constants.

The literature search completed for this version did not reveal this specific
finite-family envelope together with its equality classification,
dimension-free near-equality theorem, and scale-sharp three-cell obstruction.
We therefore formulate the contribution in terms of these concrete results,
without making a broader claim of historical priority for commutator
inequalities, almost-invariant subspaces, or stability of commutation
relations.

\section{Conclusion}

For a finite projective decomposition and a unitary transformation, the
largest aggregate cut commutator satisfies the sharp quadratic envelope
\[
  \cmax^2\le\frac{\dd^2}{2}.
\]
Positive-defect equality is rigid: it occurs exactly when two cells, and only
two cells, carry nonzero commutator defect.  Near equality inherits a
dimension-free quantitative form.  At defect scale at least $\rho$,
near-saturation with deficit $\etaq$ forces all but
\[
  2\etaq+
  \frac{36\etaq}{(1-3\etaq)\rho^2}
\]
of the relative quadratic budget onto two cells.

An explicit three-cell family proves that the inverse-square scale cannot be
eliminated: the asymptotic scaled ratio tends to $4$.  This closes the
qualitative--quantitative--sharpness chain for the finite cut envelope and
provides a machine-checked basis for further study of near-extremizers of
projective record-transfer inequalities.

\section*{Declaration of generative AI and AI-assisted technologies in the manuscript preparation process}

During the preparation of this work, the author used OpenAI ChatGPT to assist
with organization of the manuscript, language editing, literature
organization, and consistency checks between the mathematical exposition and
the Lean formalization. After using this service, the author reviewed and
edited the content as needed and takes full responsibility for the content of
the publication.

\bibliographystyle{plain}
\bibliography{references_v02}

\end{document}
